\documentclass[a4paper]{article}
% Español (México) con XeLaTeX: fontspec para fuentes Unicode, polyglossia
% para la división silábica y los nombres del documento en la variante mexicana.
\usepackage{fontspec}
\usepackage{polyglossia}
\setdefaultlanguage[variant=mexican]{spanish}
% Los nombres chinos van como «pinyin (original)»: xeCJK da a los caracteres
% Han su propia fuente; sin ella XeLaTeX los omite con un aviso.
% FandolSong no cubre algunos caracteres de nombres (U+73FA); WenQuanYi Zen Hei sí.
\usepackage[AutoFallBack=true]{xeCJK}
\setCJKmainfont{FandolSong-Regular.otf}
\setCJKfallbackfamilyfont{\CJKrmdefault}{WenQuanYi Zen Hei}
% polyglossia no traduce los nombres de listings.
\AtBeginDocument{\providecommand{\lstlistingname}{}\renewcommand{\lstlistingname}{Listado}%
  \providecommand{\lstlistlistingname}{}\renewcommand{\lstlistlistingname}{Índice de listados}}
\usepackage{amsmath, amssymb}
\usepackage{graphicx}
\usepackage[margin=2.15cm]{geometry}
\usepackage[most]{tcolorbox}
\usepackage{booktabs}
\usepackage{tabularx}
\usepackage{array}
\usepackage{float}
\usepackage{xcolor}
\usepackage{hyperref}
\usepackage{fancyhdr}
\setCJKmainfont[AutoFakeBold=2.4]{AR PL UMing CN}
\setCJKsansfont[AutoFakeBold=2.4]{Droid Sans Fallback}
\setCJKmonofont[AutoFakeBold=2.4]{Droid Sans Fallback}
\hypersetup{colorlinks=true, linkcolor=blue!55!black, urlcolor=blue!60!black}
\graphicspath{{./}{frames/}}
\setlength{\parskip}{0.55em}
\setlength{\parindent}{2em}
\setlength{\emergencystretch}{3em}
\renewcommand{\arraystretch}{1.18}
\newtcolorbox{knowledgebox}[1]{enhanced, colback=blue!5!white, colframe=blue!70!black, colbacktitle=blue!70!black, coltitle=white, fonttitle=\bfseries, title={#1}, sharp corners, breakable}
\newtcolorbox{importantbox}[1]{enhanced, colback=yellow!10!white, colframe=yellow!75!black, colbacktitle=yellow!75!black, coltitle=black, fonttitle=\bfseries, title={#1}, sharp corners, breakable}
\newtcolorbox{warningbox}[1]{enhanced, colback=red!5!white, colframe=red!70!black, colbacktitle=red!70!black, coltitle=white, fonttitle=\bfseries, title={#1}, sharp corners, breakable}
\pagestyle{fancy}
\fancyhf{}
\fancyfoot[C]{\thepage}
\fancyfoot[R]{\footnotesize\color{gray}\href{https://github.com/hqhq1025/ai-course-notes}{AI Course Notes}}
\renewcommand{\headrulewidth}{0pt}
\renewcommand{\footrulewidth}{0pt}
\newcommand{\videourl}{https://www.youtube.com/watch?v=eiQFomOuCJs}
\newcommand{\repourl}{https://github.com/hqhq1025/ai-course-notes}

\begin{document}
\begin{titlepage}
\centering
{\Large Notas de una clase técnica con pantalla compartida\par}
\vspace{1.0cm}
{\huge\bfseries Modelos fundacionales para robots y VLA: hoja de ruta de los artículos clásicos}\par
\vspace{0.5cm}
{\Large Zhang Xiaojun Podcast conversa con Chen Jianyu (陈建宇)\par}
\vspace{0.35cm}
{\large Elaboradas a partir del video con pantalla compartida, la transcripción, la estructura de capítulos, la descripción del video y los fotogramas clave de la pantalla compartida\par}
\vspace{0.3cm}
{\large 2025-04-07\par}
\vspace{0.85cm}
\includegraphics[width=0.88\textwidth,height=0.42\textheight,keepaspectratio]{cover.jpg}\par
\vfill
\begin{tcolorbox}[width=0.92\textwidth, colback=black!2!white, colframe=black!55, sharp corners]
\textbf{Título del video}: Análisis, artículo por artículo, de los modelos fundacionales para robots y los artículos clásicos de VLA (con versión de pantalla compartida): «el ser humano es el VLA más inteligente»\par
\textbf{Canal del video}: Zhang Xiaojun Podcast\par
\textbf{Duración del video}: 02:29:41\par
\textbf{Enlace del video}: \href{\videourl}{\nolinkurl{\videourl}}\par
\textbf{Nota sobre el material}: YouTube no ofrece subtítulos descargables; el texto se elaboró a partir de una transcripción por bloques con faster-whisper large-v3 (CPU, int8), la descripción del video, la estructura de capítulos y los fotogramas clave de la pantalla compartida. Este video contiene mucho contenido en pantalla compartida, por lo que el texto usa fotogramas reales de esa pantalla como eje visual.\par
\textbf{Página del proyecto}: \href{\repourl}{\nolinkurl{\repourl}}\par
\end{tcolorbox}
\end{titlepage}

\begingroup
\small
\setlength{\parskip}{0pt}
\renewcommand{\baselinestretch}{0.92}\selectfont
\tableofcontents
\endgroup
\newpage

\section{Introducción: por qué los robots necesitan un modelo base}

Esta sección plantea primero el problema de todo el episodio. El juicio central de Chen Jianyu (陈建宇) es que, en el pasado, la robótica solía funcionar con «100 escenarios, 100 tareas, 100 modelos especializados», un enfoque que difícilmente escala. El éxito de los modelos de lenguaje de gran tamaño y de los modelos multimodales le mostró a la robótica otro camino: ¿se puede entrenar para los robots un foundation model más general, que generalice entre tareas, entre cuerpos de robot y entre escenarios?

\begin{importantbox}{Tesis central de este episodio}
VLA, es decir, Vision-Language-Action Model, no consiste solo en conectar un modelo de lenguaje a un robot, sino en integrar visión, lenguaje y acción en un mismo sistema de extremo a extremo y escalable. Un verdadero modelo base para robots debe pasar de «usar los modelos fundacionales existentes para hacer robótica» a «preentrenar modelos fundacionales para los robots».
\end{importantbox}

\begin{figure}[H]
\centering
\includegraphics[width=0.96\textwidth]{ch03-00693.jpg}
\caption{How robotic learning works now? El aprendizaje robótico tradicional suele organizarse en torno al ciclo cerrado de tarea, datos, modelo y ejecución. Fotograma de la proyección, aprox. 00:11:33.}
\end{figure}

\begin{knowledgebox}{Lectura de la figura: la presión para pasar de modelos especializados a modelos generales}
La figura no presenta un artículo concreto, sino el contexto de toda la clase: el aprendizaje robótico ha dependido durante mucho tiempo de tareas específicas, datos específicos y despliegues específicos. Chen Jianyu subraya que este enfoque no es lo bastante general para hogares, fábricas y robots humanoides, por lo que se necesita la vía de los modelos base.
\end{knowledgebox}

\subsection{Las dos etapas de la IA para robots}

Esta sección distingue con claridad las dos vías. La primera etapa es leveraging foundation models in robotics, es decir, tomar modelos fundacionales existentes, como LLM, VLM o Code LM, para sustituir algún módulo del sistema robótico. La segunda etapa es pretraining foundation models for robotics, es decir, entrenar directamente para los robots modelos capaces de producir acciones. La primera consiste en ensamblar y apoyarse en lo existente; solo la segunda constituye un verdadero modelo base para robots.

\begin{figure}[H]
\centering
\includegraphics[width=0.90\textwidth]{ch05-01297.jpg}
\caption{Foundation models for robotics: relación entre módulos, desde LLM y VLA hasta Actuation. Fotograma de la proyección, aprox. 00:21:37.}
\end{figure}

\begin{knowledgebox}{Lectura de la figura: los tres bloques del robot y su sustitución por modelos fundacionales}
La figura reúne Perception, Actuation y LLM. Tradicionalmente, un robot se divide en planificación, percepción y ejecución; la ola de los modelos grandes sustituyó primero a Planning, después a Perception, y solo al final intentó que Action también se integrara en un modelo unificado.
\end{knowledgebox}

\begin{knowledgebox}{Términos clave: VLA, VLM, LLM}
\noindent\begin{tabularx}{\textwidth}{p{0.18\textwidth}p{0.34\textwidth}X}
\toprule
Término & Significado & Función en esta clase \\
\midrule
LLM & Large Language Model, modelo de lenguaje de gran tamaño & Primero sustituye o refuerza la planificación del robot. \\
VLM & Vision-Language Model, modelo de visión y lenguaje & Se usa para la percepción, la comprensión de la escena y la retroalimentación visual. \\
VLA & Vision-Language-Action Model, modelo de visión, lenguaje y acción & Conecta directamente visión, lenguaje y acción en un control de extremo a extremo. \\
Robot Foundation Model & Modelo base para robots & Su objetivo es generalizar entre tareas, entre cuerpos de robot y entre escenarios. \\
\bottomrule
\end{tabularx}
\end{knowledgebox}

\subsection{Resumen de la sección}

La aparición de VLA es la pista fundamental del paso de la robótica desde el ensamblaje de módulos hacia los modelos base generales. El hilo conductor de toda la clase va de apoyarse en LLM/VLM, a RT-2/OpenVLA/HiRT/pi0/GR00T N1, y de ahí a diffusion policy, RDT y los futuros modelos unificados de comprensión y predicción.
\section{Primera etapa: adaptar los tres bloques de la robótica con los modelos fundacionales existentes}

La sección anterior planteó una ruta de dos etapas; esta sección examina la primera: incorporar los LLM, VLM y Code LM existentes al sistema robótico. Los tres bloques de la robótica tradicional son Planning, Perception y Actuation. Los modelos grandes sustituyen primero a la planificación, porque los modelos de lenguaje son buenos para descomponer una tarea en pasos; después, los VLM sustituyen a la percepción, porque pueden vincular la escena visual con la tarea expresada en lenguaje; y, un paso más allá, los Code LM intentan automatizar también la ejecución.

\subsection{SayCan: lo que puedo hacer, no lo que digo que haga}

El problema central de SayCan es este: el modelo de lenguaje propone muchas acciones, pero el robot no necesariamente puede realizarlas en el entorno actual. SayCan combina la descomposición de tareas del modelo de lenguaje con la affordance del robot, de modo que considera tanto «qué debería hacerse según el lenguaje» como «si el robot realmente puede hacerlo». Es un trabajo pionero en la conexión de los LLM con la planificación robótica.

\begin{figure}[H]
\centering
\includegraphics[width=0.96\textwidth]{ch07-01537.jpg}
\caption{SayCan: Query LLM to rank action primitives, y después se usa la affordance function para juzgar la viabilidad. Fotograma de la pantalla proyectada, aprox. 00:25:37.}
\end{figure}

\begin{knowledgebox}{Lectura de la figura: el LLM aporta el objetivo; la affordance, la restricción}
A la izquierda de la figura está la clasificación de acciones que hace el modelo de lenguaje; a la derecha, la evaluación de las posibilidades de acción del robot en el entorno real. Al leer esta figura conviene notar que SayCan no deja que el modelo de lenguaje controle directamente al robot, sino que añade una restricción entre el plan expresado en lenguaje y la viabilidad física.
\end{knowledgebox}

\begin{warningbox}{Un plan en lenguaje no equivale directamente a una acción del robot}
El robot debe considerar su capacidad de agarre, la posición en el espacio, el estado de los objetos y los fallos de ejecución. Si el LLM solo emite «toma la taza», no sabe si la taza está al alcance, si la pinza es adecuada ni si el entorno lo permite. El valor de SayCan está en aterrizar la intención expresada en lenguaje en la affordance.
\end{warningbox}

\subsection{Inner Monologue, DoReMi y VoxPoser}

Este apartado examina tres trabajos que siguen aplicando modelos fundacionales al razonamiento robótico. Inner Monologue subraya que, durante la ejecución, el robot necesita retroalimentación verbalizada y un monólogo interior que reincorpore la retroalimentación del entorno a la planificación. DoReMi se ocupa de cómo detectar y recuperarse cuando el plan y la ejecución no coinciden. VoxPoser, por su parte, usa el modelo de lenguaje para construir 3D value maps y convierte los objetivos expresados en lenguaje en campos de valor espaciales componibles, aplicados a la manipulation.

\begin{figure}[H]
\centering
\includegraphics[width=0.90\textwidth]{ch08-01649.jpg}
\caption{Inner Monologue: un ciclo de razonamiento corporizado formado por la planificación del modelo de lenguaje y la retroalimentación del entorno. Fotograma de la pantalla proyectada, aprox. 00:27:29.}
\end{figure}

\begin{knowledgebox}{Lectura de la figura: el robot necesita narrarse a sí mismo lo que ocurrió tras ejecutar}
La figura de Inner Monologue une Action, Environmental Feedback y Corrective Action en un ciclo cerrado. No se trata simplemente de que el robot «sepa hablar», sino de que el resultado de la ejecución regrese al proceso de planificación en lenguaje y forme un estado intermedio corregible.
\end{knowledgebox}

\begin{figure}[H]
\centering
\includegraphics[width=0.94\textwidth]{ch09-01737.jpg}
\caption{DoReMi: detectar y recuperarse de la plan-execution misalignment. Fotograma de la pantalla proyectada, aprox. 00:28:57.}
\end{figure}

\begin{knowledgebox}{Lectura de la figura: la detección de desajustes es un problema real de la inteligencia corporizada}
DoReMi trata el desajuste entre planificación y ejecución: el plan dice A, pero durante la ejecución el entorno se convierte en B, o el robot no completa la tarea como se esperaba. Nos recuerda que un sistema robótico no termina al generar un plan una sola vez, sino que debe detectar, corregir y recuperarse de forma continua.
\end{knowledgebox}

\begin{figure}[H]
\centering
\includegraphics[width=0.94\textwidth]{ch11-01956.jpg}
\caption{VoxPoser: usar el modelo de lenguaje para componer 3D value maps que guían la manipulación del robot. Fotograma de la pantalla proyectada, aprox. 00:32:36.}
\end{figure}

\begin{knowledgebox}{Lectura de la figura: del objetivo en lenguaje al campo de manipulación 3D}
La figura de VoxPoser convierte las instrucciones en lenguaje en un value map en el espacio. Al leerla hay que fijarse en la palabra «componible»: distintas restricciones pueden superponerse para formar un solo objetivo de manipulación, por ejemplo, acercarse al objeto, evitar obstáculos o mantener la postura.
\end{knowledgebox}

\subsection{De los Code LM a un verdadero modelo robótico}

La idea de que los Code LM sustituyan a Actuation consiste en que el modelo escriba código para el robot o llame a las API de bajo nivel. Sin embargo, Chen Jianyu (陈建宇) señala que esto todavía no es un robot foundation model completo. Se trata más bien de orquestar herramientas y controladores existentes; un verdadero VLA requiere que el propio modelo sea capaz de emitir acciones directamente a partir de la visión y el lenguaje, y no solo de escribir código de llamadas.

\begin{importantbox}{Los límites de la primera etapa}
Leveraging foundation models permite reforzar con rapidez la planificación, la percepción y la llamada a herramientas, pero en esencia sigue siendo un sistema modular. Un robot de propósito general necesita, en última instancia, entrenar con datos de robots un modelo capaz de manejar action.
\end{importantbox}

\subsection{Resumen de la sección}

La aportación de la primera etapa es demostrar que los LLM/VLM pueden ayudar al robot a comprender tareas, planificar acciones, detectar desajustes y construir objetivos espaciales. Sin embargo, estos métodos siguen dependiendo de los módulos robóticos tradicionales y todavía no constituyen un VLA de extremo a extremo.
\section{Segunda etapa: VLA y modelos base para robots}

La sección anterior explicó cómo los modelos fundacionales existentes llegan a los robots; esta sección entra en la segunda etapa: preentrenar directamente un foundation model para robots. El objetivo de VLA es reunir visión, lenguaje y acción en un solo modelo, de modo que el modelo produzca acciones del robot directamente a partir de observaciones e instrucciones. Chen Jianyu (陈建宇) lo explica con la idea de que «el ser humano es un VLA Agent muy inteligente»: el ser humano ve el mundo, entiende el lenguaje y produce acciones; en sí mismo es un VLA general.

\begin{figure}[H]
\centering
\includegraphics[width=0.92\textwidth]{ch13-02358.jpg}
\caption{VLA Model: asigna las entradas de visión y lenguaje a una salida de action. Fotograma de la pantalla proyectada, aprox. 00:39:18.}
\end{figure}

\begin{knowledgebox}{Lectura de la figura: la A de VLA es la verdadera dificultad}
La figura va de vision y language a action y muestra la forma mínima de un VLA. V y L pueden aprovechar los avances de los VLM/LLM, pero A tiene que resolver la frecuencia, el control continuo, las trayectorias, la estabilidad y las diferencias entre cuerpos robóticos.
\end{knowledgebox}

\subsection{ALOHA y Mobile ALOHA: hardware de bajo costo y datos de demostración}

Esta subsección vuelve del concepto de VLA a la base de datos y hardware. La razón es que un modelo base para robots no surge solo de la estructura de un modelo grande descrita en un artículo; antes que nada necesita un sistema barato, replicable y capaz de recolectar demostraciones de forma estable. Aquí ALOHA cumple el papel de «producir datos reales de manipulación con dos manos», y ofrece una referencia concreta para la discusión posterior sobre action chunk, aprendizaje por imitación y generalización entre tareas.

La serie ALOHA no es un VLA moderno en sentido estricto, pero es importante para el aprendizaje robótico: hardware de dos brazos de bajo costo, demostraciones por teleoperation, Action Chunking Transformer, y la capacidad de que el robot imite manipulaciones finas con dos manos. Mobile ALOHA agrega una base móvil, extiende la manipulación con dos brazos a escenarios móviles más complejos y se dio a conocer ampliamente gracias a una gran cantidad de videos de demostración.

\begin{figure}[H]
\centering
\includegraphics[width=0.96\textwidth]{ch15-02658.jpg}
\caption{ALOHA: Action Chunking Transformer para predecir secuencias de acciones en manipulación con dos manos. Fotograma de la pantalla proyectada, aprox. 00:44:18.}
\end{figure}

\begin{knowledgebox}{Lectura de la figura: Action Chunking es predicción de secuencias de acciones}
La figura muestra la estructura de ACT: recibe observaciones y produce una secuencia de acciones futuras. Lo esencial es el chunking: el modelo no predice solo la siguiente acción, sino un tramo de trayectoria de acciones, que luego se ejecuta en horizonte deslizante. La intuición es parecida a la del control predictivo basado en modelo.
\end{knowledgebox}

\begin{figure}[H]
\centering
\includegraphics[width=0.92\textwidth]{ch16-02943.jpg}
\caption{Mobile ALOHA: extiende la manipulación de bajo costo con dos brazos a la teleoperación móvil de cuerpo completo. Fotograma de la pantalla proyectada, aprox. 00:49:03.}
\end{figure}

\begin{warningbox}{Limitaciones de ALOHA}
ALOHA es muy potente, pero se parece más a un sistema de aprendizaje por imitación de alta calidad. Sigue dependiendo en buena medida de la tarea y de los datos, y no equivale a un modelo base VLA capaz de generalizar entre una gran cantidad de robots y tareas.
\end{warningbox}

\subsection{Gato, RT-1 y Octo}

El siguiente grupo de trabajos lleva la pregunta de «¿se pueden recolectar demostraciones de alta calidad?» a «¿se puede entrenar una política más general?». De Gato a RT-1 y luego a Octo, el cambio central es que la idea del agente general se va concretando en datos reales de robots y en marcos de policy de código abierto. El docente enfatiza que no todos estos trabajos son VLA en el sentido actual, pero en conjunto abrieron el camino hacia el modelado unificado, el entrenamiento multitarea y el escalamiento del control robótico.

Gato es un intento temprano de agente general en el que «un solo modelo atiende muchas tareas»; la idea fue muy adelantada, pero en ese momento los modelos, los datos y la capacidad de cómputo aún no estaban suficientemente maduros. RT-1, en cambio, se orienta de forma más directa al control robótico y entrena un Robotics Transformer con datos de control del mundo real. Octo es una política robótica general de código abierto que pone el acento en reunir distintas tareas y datos en un marco de policy reutilizable.

\begin{figure}[H]
\centering
\includegraphics[width=0.92\textwidth]{ch17-03097.jpg}
\caption{Gato: A Generalist Agent, que convierte en tokens múltiples tareas y las modela de forma unificada. Fotograma de la pantalla proyectada, aprox. 00:51:37.}
\end{figure}

\begin{knowledgebox}{Lectura de la figura: Gato ya se acercaba conceptualmente a VLA}
La figura de Gato muestra entradas multitarea y multimodales unificadas en un solo modelo. Su valor histórico es haber propuesto la dirección del generalist agent, aunque todavía no contaba con la escala de modelo, los datos ni la capacidad de ejecución robótica de los VLM/VLA actuales.
\end{knowledgebox}

\begin{figure}[H]
\centering
\includegraphics[width=0.94\textwidth]{ch18-03372.jpg}
\caption{RT-1: Robotics Transformer for real-world control at scale. Fotograma de la pantalla proyectada, aprox. 00:56:12.}
\end{figure}

\begin{knowledgebox}{Lectura de la figura: RT-1 es un hito en el escalamiento de datos robóticos}
RT-1 introduce en un Transformer la descripción de la tarea, las observaciones de imagen y los token de acción. Su aporte no se limita a la estructura del modelo: también incluye los datos de control de robots reales y un flujo de entrenamiento escalable.
\end{knowledgebox}

\begin{figure}[H]
\centering
\includegraphics[width=0.92\textwidth]{ch19-03650.jpg}
\caption{Octo: política robótica general de código abierto que permite convertir en tokens la task y la observation. Fotograma de la pantalla proyectada, aprox. 01:00:50.}
\end{figure}

\subsection{Aprendizaje entre cuerpos robóticos: CrossFormer y la serie GR}

Un modelo robótico general no puede dominar solo un tipo de brazo mecánico o una sola tarea. CrossFormer se enfoca en el cross-embodied learning y busca que una sola policy cubra cuerpos robóticos distintos, como manipulation, navigation, locomotion y aviation. Por su parte, GR-1/GR-2, del AI Lab de ByteDance (字节), incorporan a la manipulación robótica el preentrenamiento con generación de video, video-language-action y web-scale knowledge, y ponen el acento en aprender de videos y conocimiento a gran escala.

\begin{figure}[H]
\centering
\includegraphics[width=0.94\textwidth]{ch20-03892.jpg}
\caption{CrossFormer: Scaling cross-embodied learning, una sola policy orientada a varios tipos de cuerpos robóticos. Fotograma de la pantalla proyectada, aprox. 01:04:52.}
\end{figure}

\begin{knowledgebox}{Lectura de la figura: la transferencia entre cuerpos robóticos es esencial para la generalización de VLA}
La figura de CrossFormer destaca una política compartida entre distintos cuerpos robóticos. Un modelo base para robots verdaderamente general no puede servir a un único hardware; de lo contrario, se vuelve a «un modelo por robot».
\end{knowledgebox}

\begin{figure}[H]
\centering
\includegraphics[width=0.94\textwidth]{ch21-04284.jpg}
\caption{GR-1: trabajo de manipulación robótica visual basado en preentrenamiento con generación de video. Fotograma de la pantalla proyectada, aprox. 01:11:24.}
\end{figure}

\subsection{Resumen de la sección}

La segunda etapa empieza a poner en el centro los datos robóticos, la salida de acciones y la generalización entre cuerpos robóticos. ALOHA aporta demostraciones de bajo costo y aprendizaje de secuencias de acciones; trabajos como RT-1/Octo/RT-X/OpenVLA llevan los VLA hacia políticas robóticas más generales.
\section{RT-2, RT-X, OpenVLA y Action Head}

Ya se explicó que los modelos base para robots necesitan datos de robots reales y políticas que funcionen entre tareas; esta sección avanza hacia el grupo de artículos más cercano a la línea principal actual de VLA. Aquí la pregunta cambia: si un VLM ya aprendió una gran cantidad de conocimiento visual y lingüístico a partir de imágenes y textos de internet, ¿puede el robot transferir ese conocimiento a la salida de acciones, en lugar de volver a aprender el sentido común sobre el mundo a partir de una pequeña cantidad de datos de robots?

Esta sección entra en la línea principal actual de VLA. RT-2 se considera uno de los trabajos pioneros de VLA: transfiere el conocimiento visual y lingüístico a escala web al control de robots. La idea central es esta: un VLM preentrenado ya domina una gran cantidad de conocimiento visual y lingüístico; ¿es posible, mediante datos de acciones de robots, hacer que produzca directamente acciones (action)? Esto conecta el conocimiento de internet con el control de robots.

\begin{figure}[H]
\centering
\includegraphics[width=0.96\textwidth]{ch23-05001.jpg}
\caption{RT-2: Vision-Language-Action Models transfer web knowledge to robotic control. Fotograma de la pantalla proyectada, aprox. 01:23:21.}
\end{figure}

\begin{knowledgebox}{Lectura de la figura: RT-2 usa un VLM como backbone para producir acciones del robot}
En la figura se observa que RT-2 conecta el VLM con las acciones (action) del robot. El núcleo es la transferencia: el modelo no aprende desde cero solo con datos de robots, sino que transfiere el conocimiento visual y lingüístico de internet al control de robots.
\end{knowledgebox}

\subsection{RT-X y Open X-Embodiment}

Esta subsección desplaza el foco de la estructura de un modelo individual a la infraestructura de datos. RT-2 demostró que transferir un VLM al control de robots tiene valor, pero si los datos de entrenamiento provienen solo de unos pocos robots y unas pocas tareas, el modelo sigue quedando fácilmente atado a la plataforma robótica, al escenario y al espacio de acciones. Por eso, la pregunta central de RT-X/Open X-Embodiment es cómo reunir la experiencia dispersa en distintos laboratorios y distintos robots en datos generales con los que se pueda entrenar.

RT-X y Open X-Embodiment ponen el énfasis en el conjunto de datos y en los modelos que funcionan entre plataformas robóticas. Los datos de robots son escasos y fragmentados, y a un solo laboratorio le resulta difícil cubrir suficientes tareas y plataformas robóticas. Open X-Embodiment reúne datos de múltiples robots y múltiples tareas para intentar entrenar modelos que funcionen entre plataformas robóticas, y muestra la posibilidad de que una política general supere a las políticas especializadas en una sola tarea.

\begin{figure}[H]
\centering
\includegraphics[width=0.96\textwidth]{ch24-05336.jpg}
\caption{RT-X / Open X-Embodiment: conjunto de datos de múltiples robots y el modelo RT-X. Fotograma de la pantalla proyectada, aprox. 01:28:56.}
\end{figure}

\begin{knowledgebox}{Lectura de la figura: el conjunto de datos es en sí mismo infraestructura}
La figura reúne una gran cantidad de robots y escenarios de tareas. Al leerla, conviene notar que el cuello de botella de VLA no es solo la estructura del modelo, sino también la estandarización de datos entre laboratorios, entre plataformas de hardware y entre tareas.
\end{knowledgebox}

\subsection{OpenVLA y HiRT: código abierto y procesamiento jerárquico de acciones}

OpenVLA se aproxima a una versión de código abierto de RT-2: conecta un VLM a la salida de acciones y publica como código abierto el modelo y el proceso de entrenamiento. HiRT, por su parte, señala que producir acciones directamente con un VLM presenta problemas de frecuencia de acción y de control fino, por lo que añade una policy/head dedicada a procesar la acción (Action): el VLM se encarga, a baja frecuencia, de la comprensión y la planificación, y la Action Policy ejecuta el control a alta frecuencia.

\begin{figure}[H]
\centering
\includegraphics[width=0.96\textwidth]{ch25-05531.jpg}
\caption{OpenVLA: modelo VLA de código abierto que conecta un VLM con la salida de acciones del robot. Fotograma de la pantalla proyectada, aprox. 01:32:11.}
\end{figure}

\begin{figure}[H]
\centering
\includegraphics[width=0.96\textwidth]{ch26-05765.jpg}
\caption{HiRT: mejora del control de robots mediante Robot Transformers jerárquicos. Fotograma de la pantalla proyectada, aprox. 01:36:05.}
\end{figure}

\begin{knowledgebox}{Lectura de la figura: el Action Head es un complemento importante para llevar VLA a la ingeniería}
En la figura de HiRT, la información del VLM de alto nivel se transmite a la Action Policy. Esto corresponde a un hecho de ingeniería: un VLM grande infiere con lentitud, el control de robots opera a alta frecuencia, y la salida de acciones necesita un módulo dedicado para manejar el control continuo y la retroalimentación visual local.
\end{knowledgebox}

\subsection{Figure Helix, pi0 y GR00T N1}

El siguiente tramo entra en la industria y en sistemas más cercanos a la conversión en producto. OpenVLA y HiRT ya dejaron al descubierto una tensión central: el VLM es bueno para la comprensión a baja frecuencia y la planificación semántica, pero un robot real necesita acciones (action) de alta frecuencia, continuas y estables. Lo que Figure Helix, pi0 y GR00T N1 tienen en común es que separan con mayor claridad el «módulo de comprensión» y el «módulo de acción», y luego los vuelven a conectar mediante el diseño del sistema.

Aunque Figure Helix no ha publicado un artículo, representa la arquitectura más reciente de la industria: el Sistema 2, en el nivel superior, funciona como un VLM preentrenado, y el Sistema 1, en el nivel inferior, funciona como control de acciones en tiempo real. pi0 introduce un flow model en VLA y usa un Action Expert para el control general de robots. NVIDIA GR00T N1, por su parte, pone el énfasis en un modelo fundacional abierto para robots humanoides y también adopta una estructura de VLM más procesamiento de acciones.

\begin{figure}[H]
\centering
\includegraphics[width=0.92\textwidth]{ch27-05946.jpg}
\caption{Figure AI Helix: esquema de la arquitectura de Sistema 2 y Sistema 1. Fotograma de la pantalla proyectada, aprox. 01:39:06.}
\end{figure}

\begin{figure}[H]
\centering
\includegraphics[width=0.92\textwidth]{ch28-06038.jpg}
\caption{pi0: Vision-Language-Action Flow Model for General Robot Control. Fotograma de la pantalla proyectada, aprox. 01:40:38.}
\end{figure}

\begin{figure}[H]
\centering
\includegraphics[width=0.94\textwidth]{ch29-06126.jpg}
\caption{NVIDIA GR00T N1: modelo fundacional abierto para robots humanoides de propósito general. Fotograma de la pantalla proyectada, aprox. 01:42:06.}
\end{figure}

\subsection{Resumen de la sección}

RT-2 transfiere el conocimiento de un VLM al control de robots; RT-X/Open X-Embodiment pone el énfasis en la escala de los datos y en el funcionamiento entre plataformas robóticas; OpenVLA ofrece una vía de código abierto, y HiRT/pi0/GR00T N1 comienzan a tratar con más seriedad el módulo de acción (Action). Esto indica que la dificultad central de VLA se está desplazando de «conectar el lenguaje y la visión» a «cómo producir acciones con alta frecuencia, de forma estable y con capacidad de generalización».
\section{Diffusion Policy, RDT y la línea de los modelos del mundo}

Los capítulos anteriores giraron principalmente en torno a Transformer/VLM/VLA; este capítulo examina la línea de generación de action. Diffusion Policy aplica los modelos de difusión al aprendizaje de políticas visomotoras: no genera imágenes, sino trayectorias de acciones. Esto resulta natural para los robots, porque una secuencia de acciones también puede verse como un objeto que debe eliminarse de ruido y generarse paso a paso.

\begin{figure}[H]
\centering
\includegraphics[width=0.96\textwidth]{ch30-06320.jpg}
\caption{Diffusion Policy: aprendizaje de una visuomotor policy mediante difusión de acciones. Fotograma de la pantalla proyectada, aprox. 01:45:20.}
\end{figure}

\begin{knowledgebox}{Lectura de la figura: los modelos de difusión también pueden generar action}
La figura muestra los procesos de entrenamiento y de inferencia de Diffusion Policy. La transferencia central es esta: los modelos de difusión no solo sirven para generar imágenes, sino también para generar action trajectory, y son especialmente adecuados para acciones continuas y trayectorias multimodales.
\end{knowledgebox}

\subsection{RDT-1B: un modelo fundacional de difusión para manipulación bimanual}

Esta sección sigue indagando cómo debe modelarse la action misma. La línea VLM+Action Head trata el módulo de acción como un controlador de alta frecuencia, pero aún hace falta un mecanismo generativo más adecuado para trayectorias continuas, elecciones multimodales y coordinación de ambas manos. La importancia de RDT-1B consiste en llevar la intuición de Diffusion Policy a la escala de un modelo fundacional, de modo que la «generación de acciones» deje de ser solo un componente accesorio de una pequeña red de política.

RDT-1B extiende Diffusion Policy a un modelo fundacional más grande para manipulación bimanual. Busca resolver el unified action space, las distintas action head y la manipulation con dos brazos. Su importancia reside en llevar la difusión de acciones propia de los modelos pequeños hacia una escala mayor, más cercana a un foundation model.

\begin{figure}[H]
\centering
\includegraphics[width=0.96\textwidth]{ch31-06571.jpg}
\caption{RDT-1B: un modelo fundacional de difusión para robots de manipulación bimanual. Fotograma de la pantalla proyectada, aprox. 01:49:31.}
\end{figure}

\subsection{Prediction with Action y VPP}

Prediction with Action coloca la predicción visual y la predicción de acciones en un proceso conjunto de eliminación de ruido, con el fin de conectar los modelos del mundo con los VLA. Su continuación, VPP, pone el énfasis en usar predictive visual representations para formar una generalist robot policy. Lo central aquí no es solo producir la action, sino comprender a la vez los cambios visuales futuros y las consecuencias de las acciones.

\begin{figure}[H]
\centering
\includegraphics[width=0.96\textwidth]{ch32-07061.jpg}
\caption{Prediction with Action / VPP: aprendizaje de políticas visuales mediante eliminación de ruido conjunta. Fotograma de la pantalla proyectada, aprox. 01:57:41.}
\end{figure}

\begin{knowledgebox}{Lectura de la figura: los modelos del mundo entran en los VLA}
La figura incluye la predicción de video, las acciones y el proceso de eliminación de ruido. Al leerla hay que retener una idea: el robot no solo necesita saber cuál es la siguiente acción, sino también predecir cómo esa acción cambiará el mundo visual; esa es la línea de los modelos del mundo.
\end{knowledgebox}

\subsection{Direcciones futuras: UP-VLA y aprendizaje por refuerzo en línea}

Esta sección pasa de los sistemas existentes a los problemas de la siguiente etapa. Los VLA actuales ya pueden conectar visión, lenguaje y acción, pero todavía no resuelven de forma estable la unificación entre «comprender el mundo» y «predecir las consecuencias de las acciones», ni el ciclo cerrado de mejora continua en entornos reales. Por ello, Chen Jianyu (陈建宇) condensa las direcciones futuras en dos líneas: una estructura de modelo más unificada y un aprendizaje en línea más seguro y controlable.

Por último, Chen Jianyu enumera dos direcciones futuras. UP-VLA busca unificar understanding y prediction, integrando la comprensión y la predicción en un solo modelo del embodied agent. La otra dirección es el online reinforcement learning, que usa la retroalimentación en línea para seguir mejorando el VLA. La dificultad es que, si se aplica RL directamente a todo el VLM/VLA, el entrenamiento puede volverse inestable; una estrategia consiste en congelar primero el VLM, entrenar solo la action head y después liberarlo de manera gradual.

\begin{figure}[H]
\centering
\includegraphics[width=0.96\textwidth]{ch33-07592.jpg}
\caption{iRe-VLA / Online RL: mejora de un Vision-Language-Action Model mediante aprendizaje por refuerzo en línea. Fotograma de la pantalla proyectada, aprox. 02:06:32.}
\end{figure}

\begin{importantbox}{Dos líneas futuras}
La primera línea es unificar comprensión, predicción y acción, de modo que el modelo no solo entienda el mundo, sino que también pueda predecir las consecuencias de sus acciones; la segunda línea es el aprendizaje por refuerzo en línea, que permite al robot mejorar continuamente a partir de la retroalimentación de entornos reales o simulados.
\end{importantbox}

\subsection{Resumen de la sección}

Diffusion Policy y RDT llevan el modelado de la generación de acciones hacia un control continuo más potente; Prediction with Action/VPP introduce los modelos del mundo en los VLA; UP-VLA y el RL en línea apuntan hacia los modelos unificados y el aprendizaje continuo del futuro.
\section{Síntesis y ampliación}

Esta sección condensa el episodio completo en seis conclusiones. Primera: lo decisivo en la revolución robótica es el paso de los modelos especializados a los modelos generales. Segunda: la vía inicial consistió en usar LLM/VLM/Code LM para transformar Planning, Perception y Actuation. Tercera: un VLA propiamente dicho integra Vision, Language y Action en un modelo de extremo a extremo. Cuarta: las líneas como RT-2/OpenVLA transfieren el conocimiento de internet mediante un VLM, pero el control de action a alta frecuencia todavía requiere un tratamiento específico. Quinta: Diffusion Policy/RDT muestran que la propia acción puede modelarse con modelos generativos. Sexta: las direcciones futuras son los modelos del mundo, la unificación de comprensión y predicción, y el aprendizaje por refuerzo en línea.

\begin{importantbox}{Ubicar la clase de VLA con pantalla compartida dentro de la serie de IA/internet de Zhang Xiaojun (张小珺)}
Este episodio es la base técnica de las entrevistas sobre inteligencia corporizada como EP106, EP109 y EP121. No es una historia de emprendimiento, sino una hoja de ruta de los modelos base para robots, y explica por qué VLA se volvió un término central en la discusión sobre robótica de 2025.
\end{importantbox}

\subsection{Tabla de consulta rápida de artículos clásicos}

\noindent\begin{tabularx}{\textwidth}{p{0.20\textwidth}p{0.30\textwidth}X}
\toprule
Artículo/sistema & Aportación principal & Lugar en esta clase \\
\midrule
SayCan & Planificación con LLM + viabilidad por affordance & Ejemplo representativo de sustituir Planning con un modelo de lenguaje. \\
Inner Monologue & Retroalimentación del entorno y autocorrección expresada en lenguaje & Hace que la retroalimentación de la ejecución vuelva a la planificación. \\
DoReMi & Detección y recuperación de desajustes entre el plan y la ejecución & Atiende las fallas de ejecución en robots reales. \\
VoxPoser & 3D value maps & Convierte objetivos en lenguaje en campos de operación espacial. \\
ALOHA & Manipulación bimanual de bajo costo y ACT & Datos de demostración y aprendizaje de secuencias de acciones. \\
RT-1/RT-2 & Transformer para robots / transferencia VLA & Línea principal de Google Robotics. \\
RT-X/OpenVLA & Datos de múltiples cuerpos robóticos y VLA de código abierto & Datos e infraestructura de código abierto. \\
HiRT/pi0/GR00T N1 & VLM + Action Head/Expert & Tratan la action con más rigor. \\
Diffusion Policy/RDT & Difusión de acciones y modelo base de manipulación bimanual & action policy generativa. \\
UP-VLA/iRe-VLA & Unificación de comprensión y predicción, y RL en línea & Direcciones futuras. \\
\bottomrule
\end{tabularx}

\subsection{Preguntas para el seguimiento}

Por último, volvemos al marco de observación a largo plazo de esta serie de entrevistas de IA/internet de Zhang Xiaojun: VLA no es el nombre de un modelo aislado, sino un conjunto de problemas de sistema en torno a los datos, los modelos, el hardware, la adopción comercial y el entrenamiento seguro. Las siguientes preguntas sirven como lista de seguimiento para juzgar si los modelos base para robots han entrado de verdad en una etapa de scale comparable a la de los grandes modelos de lenguaje.

\begin{enumerate}
\item ¿Aparecerá la scaling law de los VLA de manera tan estable como en los modelos de lenguaje?
\item ¿Podrán los conjuntos de datos de múltiples cuerpos robóticos ser lo bastante grandes y estandarizados para sostener políticas verdaderamente generales?
\item Entre un VLM backbone que emite la action directamente y un VLM + Action Head, ¿cuál de las dos líneas es más escalable?
\item ¿Se convertirán los modelos del mundo en el factor decisivo para la generalización robótica, o bastará con la propia action policy?
\item ¿Puede el aprendizaje por refuerzo en línea ejecutarse de forma segura, eficiente y sostenible en robots reales?
\item ¿Los robots humanoides llegarán primero a la aplicación práctica gracias al cuerpo robótico y la teleoperación, o tendrán que esperar a que maduren las capacidades de los VLA?
\end{enumerate}

\subsection{Lecturas adicionales}

\begin{itemize}
\item Si le interesa la historia académica de la inteligencia corporizada, puede compararse con la entrevista a Wang He (王鹤) en EP106.
\item Si le interesan la simulación y los datos sintéticos, puede compararse con la entrevista a Xie Chen (谢晨) de Guanglun (光轮) en EP109.
\item Si le interesan los modelos del mundo multimodales y el razonamiento visual, puede compararse con la entrevista a Zhang Xiangyu (张祥雨) en EP102.
\item Si le interesan la industrialización de la robótica y las vías de entrada por hardware, puede compararse con las entrevistas a Rokid en EP104 y a Tan Jie (谭捷) en EP121.
\end{itemize}

\end{document}
